% Part: first-order-logic
% Chapter: introduction
% Section: formulas

\documentclass[../../../include/open-logic-section]{subfiles}

\begin{document}

\olfileid{fol}{int}{fml}

\section{\usetoken{P}{formula}}

پہلے درجے دی منطق دے !!{sentence}s نوں پوری ٹھیکائی نال
مقرر کرن تے !!{variable}s دی ورتوں توں پیدا ہون والے
مسئلیاں نال نمٹن لئی اسیں ایہ طریقہ ورتاں گے۔
سب توں پہلاں اسیں اظہارات دا اک \emph{وکھرا} سیٹ
تعریف کراں گے: !!{formula}s۔ ایس توں بعد
اسیں انہاں وچ !!{variable}s دے کردار نوں ویکھ
سکاں گے---تے کجھ ہور تصورات، یعنی ”آزاد“ تے
”پابند“ !!{variable}s، دے سہارے ایہ تعریف
کر سکاں گے کہ !!a{sentence} کیہ اے (یعنی ایسا
!!a{formula} جیہدے وچ آزاد !!{variable}s نہ
ہون)۔ اسیں ایہ صرف ایس لئی نہیں کردے کہ
”!!{sentence}“ دی تعریف سنبھالن جوگی ہو جائے،
بلکہ ایس لئی وی کہ معنوی تصور، یعنی پورا اترن،
دی تعریف دے طریقے وچ ایہ فرق نہایت اہم ہووے گا۔

آؤ اک سادہ پہلے درجے دی زبان لئی ”!!{formula}“
تعریف کریے، جیہدے وچ صرف اک !!{predicate}~$\Obj P$
تے اک !!{constant}~$\Obj a$ ہون، تے منطقی
علامتاں وچ صرف $\lnot$، $\land$ تے~$\lexists$
ہون۔ ساڈیاں پوری تعریفاں ایہناں توں کہیں عام
ہون گیاں: اسیں لامتناہی تعداد وچ !!{predicate}s
تے !!{constant}s دی اجازت دیاں گے۔ اصل وچ
اسیں !!{function}s نوں وی ویکھاں گے، جیہڑے
!!{constant}s تے !!{variable}s نال مل کے
”ٹرم“ بنا سکدے نیں۔ فی الحال صرف $\Obj a$
تے متغیر ساڈے ٹرم ہون گے۔ البتہ سانوں
لامتناہی !!{variable}s ضرور چاہیدے نیں۔
باقاعدہ طور تے اسیں $\Obj v_0$، $\Obj
v_1$، \dots، نوں متغیر علامتاں ورتاں گے۔

\begin{defn}
\emph{!!{formula}s} دا سیٹ~$\Frm$ ایویں تعریف ہوندا اے:
\begin{enumerate}
\item\ollabel{fmls-atom} $\Atom{\Obj P}{\Obj a}$ تے $\Atom{\Obj
  P}{\Obj v_i}$ !!{formula}s نیں ($i \in \Nat$)۔

\tagitem{prvNot}{\ollabel{fmls-not}جے $!A$ !!a{formula} اے، تاں $\lnot !A$
  وی !!{formula} اے۔}{}

\tagitem{prvAnd}{جے $!A$ تے $!B$ !!{formula}s نیں، تاں $(!A \land
  !B)$ وی !!a{formula} اے۔}{}

\tagitem{prvEx}{\ollabel{fmls-ex}جے $!A$ !!a{formula} اے تے $x$ !!a{variable}
  اے، تاں $\lexists[x][!A]$ وی !!a{formula} اے۔}{}

\tagitem{limitClause}{\ollabel{fmls-limit}ایہناں توں سوا ہور کجھ وی !!a{formula} نہیں۔}{}
\end{enumerate}
\end{defn}

\olref{fmls-atom} سانوں دسدا اے کہ
$\Atom{\Obj P}{\Obj a}$ تے
$\Atom{\Obj P}{\Obj v_i}$، ہر $i \in
\Nat$ لئی، !!{formula}s نیں۔ ایہناں نوں
\emph{ایٹمی} !!{formula}s آکھدے نیں۔
ایہ ابتدا لئی مواد دیندے نیں۔ باقی شقّاں
پہلاں بنے فارمولیاں توں نویں !!{formula}s
بنان دے طریقے دیندیاں نیں۔ مثال لئی،
\olref{fmls-not} توں سانوں
$\lnot \Atom{\Obj P}{\Obj v_2}$ اک
!!a{formula} ملدا اے، کیوں جے
$\Atom{\Obj P}{\Obj v_2}$ پہلے ہی
\olref{fmls-atom} مطابق !!a{formula}
اے۔ فیر \olref{fmls-ex} توں
$\lexists[\Obj v_2][\lnot \Atom{\Obj P}{\Obj v_2}]$
وی اک !!{formula} ملدا اے، تے ایس
طرح اگے۔ \olref{fmls-limit} دسدا
اے کہ \emph{صرف} ایس طریقے نال
بنائیاں لڑیاں ہی !!{formula}s شمار
ہوندیاں نیں۔ خاص طور تے
$\lexists[\Obj v_0][\Atom{\Obj P}{\Obj a}]$
تے $\lexists[\Obj
  v_0][\lexists[\Obj v_0][\Atom{\Obj P}{\Obj a}]]$
!!{formula}s \emph{ہوندے} نیں؛ پر
$(\lnot \Atom{\Obj P}{\Obj a})$
فارمولا نہیں، کیوں جے اوہدے
باہر وادھو قوسین نیں۔

!!{formula}s نوں ایس طریقے نال تعریف کرن
نوں \emph{استقرائی تعریف} آکھدے نیں۔
ایہ سانوں !!{formula}s بارے استقرا
والے ثبوت دا اک روپ، یعنی
\emph{ساختی استقرا}، ورت کے قضیے
ثابت کرن دی اجازت دیندا اے۔
انہاں تصورات اُتے عمومی بحث
\olref[mth][ind][idf]{sec} تے
\olref[mth][ind][sti]{sec} وچ اے؛
اگے دے ثبوتاں وچ جان توں پہلاں
تہانوں اوہ حصہ مُڑ ویکھ لینا چاہیدا اے۔
بنیادی خیال ایہ اے: جے تسی ثابت کرنا
چاہندے او کہ کوئی گل سارے !!{formula}s
لئی سچی اے، تاں سب توں پہلاں اوہ
ایٹمی !!{formula}s لئی ثابت کرو؛
فیر دکھاؤ کہ \emph{جے} اوہ کسے
!!{formula}~$!A$ (تے~$!B$) لئی سچی
اے، تاں $\lnot !A$، $(!A \land !B)$
تے $\lexists[x][!A]$ لئی \emph{وی}
سچی اے۔ مثال لئی، ایس نال ایہ گل
$\lexists[\Obj v_2][\lnot \Atom{\Obj P}{\Obj v_2}]$
لئی ثابت ہوندی اے: پہلے حصے توں
پتا لگدا اے کہ ایٹمی !!{formula}~$\Atom{\Obj
P}{\Obj v_2}$ لئی ایہ سچی اے۔
فیر دوجے حصے توں
$\lnot \Atom{\Obj P}{\Obj v_2}$
لئی سچی ہوندی اے؛ تے اوسے حصے نوں
دوبارہ لا کے
$\lexists[\Obj v_2][\lnot \Atom{\Obj P}{\Obj v_2}]$
لئی وی۔ چونکہ سارے !!{formula}s
ایٹمی !!{formula}s توں استقرائی
طریقے نال بندے نیں، ایس لئی
ایہ دلیل انہاں وچوں ہر اک لئی
چلدی اے۔

\end{document}
